\section{Case Study}
\label{sec:case-study}

We now illustrate the optimal pilot study design in a realistic setting, using a simulated case study based on education RCT data.
While the theory from \Cref{sec:large-budget} and \Cref{sec:small-budget} holds only in the large budget regime $(b\to\infty)$ and the small budget regime $(b<\overline{b})$, in this section we show that both regimes can correspond to realistic sample sizes when we calibrate our model to real data.

For this case study, we focus on the trade-off between sampling from the most promising type and sampling a representative sample, and assume constant sampling costs across school types.
Over-sampling from the most promising type is done in practice for a variety of reasons, including making the potential treatment effect easier to detect and reducing the potential for causing harm~\cite{FDACD19}.
Our model captures the increased power of a targeted sample, and allows us to evaluate the trade-off between statistical power and generalizability.

We model a hypothetical pilot study preceding a real large-scale education RCT, the National Study of Learning Mindsets (NSLM) \cite{YeagerHa19}.
This experiment focused on a large representative sample of 6,320 lower achieving 9th grade students.
It scaled up the evaluation of an intervention designed to teach students a growth mindset, which had already been evaluated in prior RCTs.
Because there were theoretical reasons to believe that the treatment effect of the intervention would be moderated by at least two school-level features, we can model the decision of a pilot designer who could have chosen to recruit a representative sample or to focus on schools where the treatment had the best chance of success.

We begin by describing the details of our model calibration.
The follow-up is fixed to be representative with a sampling standard deviation corresponding to a sample size of 6,320 students.
This means we calibrate the \emph{effective} follow-up size $n_2$ used in the model by setting $\sqrt{1 / n_2} = 0.03$ to match the reported standard error of the NSLM study~\cite{YeagerHa19}.
We can then use the ratio between the number of students in a hypothetical pilot and the number of students in the NSLM study to calibrate the effective pilot size $n_1$.
We use one-sided significance tests at level $\alpha = 0.025$.
Since we assume costs are constant across types, the pilot budget is equivalent to the number of students sampled (the $x$-axis of \Cref{fig:simulation_2_policy_comparison}).

\begin{figure}[ht]
  \centering
  \includegraphics[width=0.78\textwidth]{figures/simulation_2_policy_comparison.pdf}
  \caption{
    \textbf{Pilot RCT Impact.}
    Comparison of the small-budget and limiting large-budget optimal pilot designs for a range of sample sizes in a semi-synthetic scenario calibrated to the NSLM study.
    In this instance where costs across types are assumed constant, the budget is equivalent to a constraint on the number of students, and the small-budget optimal design is equivalent to sampling only from the schools with the highest prior mean.
    The limiting large-budget optimal design is representative sampling.
    The top panel shows the expected change in standardized effect size achieved by the small-budget and large-budget optimal designs as well as an optimal design, while the bottom panel shows each design's improvement over a baseline whose continuation decision is purely random with the same false-positive rate, as a fraction of the optimal design's improvement
  }
  \label{fig:simulation_2_policy_comparison}
\end{figure}

We model 6 types of schools based on the school-level moderators reported in \textcite{YeagerHa19}.
They describe that theory supported two school-level moderators of the treatment effect: the school's achievement level, and the supportiveness of peer norms at the school.
They divide the schools into three achievement levels: low ($<25$th), medium ($25$-$75$th), and high ($>75$th percentile), with medium achievement schools expected to benefit most from the intervention.
Similarly, they divide schools into two supportiveness levels: low ($<50$th) and high ($>50$th percentile), with high supportiveness schools expected to benefit more from the intervention.
We define our 6 school types as corresponding to each combination of achievement and supportiveness levels, and assume they occur with frequency proportional to the product of their group frequencies (since \textcite{YeagerHa19} did not report their joint frequency distribution).

We base our prior on the distribution of effects across education RCTs.
In meta-analyses of education studies it is common to report standardized effect sizes, scaling the treatment effect to the standard deviation of the outcome~\cite{Kraft20}.
\textcite{YeagerHa19} discuss that a standardized effect size of 0.20 would be on the large end of possible effects, so we set the prior mean for the most promising type---medium achievement, high peer supportiveness---to 0.10 as a reasonable expectation of a good treatment effect before the pilot stage.
We use a multivariate normal prior, and set the prior mean for types that have one lower-prior moderator to 0.05, and types that have both lower-prior moderators to 0.025.
We set the prior standard deviation for each type to 0.28, matching the empirical value reported in a review of 747 preK--12 education RCTs \cite{Kraft20}, and we fix a common correlation value of $0.5$ between all types to represent a moderate level of heterogeneity.

\Cref{fig:simulation_2_policy_comparison} visualizes the pilot impact in terms of standardized effect size after calibrating our model as described above.
For pilot samples below 323 students, sampling from only the most promising schools is optimal (in this instance, the small-budget optimal design is equivalent to sampling only the school-type with the highest prior mean).
For pilot samples above 823 students representative sampling outperforms the small-budget optimal design, and above about 2000 students, representative sampling is close to optimal.
At least some of the RCTs that evaluated the growth mindset intervention prior to the NSLM study had enrolled ``thousands of adolescents''~\cite{YeagerHa19}, so both ends of this range of pilot sample sizes are plausible.
